\documentclass{article}
\usepackage[UTF8]{ctex}
\usepackage{amsmath,amssymb,amsfonts}
\usepackage{enumitem}
\usepackage{geometry}
\geometry{a4paper, margin=1in}
% === 错题本排版优化（不参与题目标识） ===
\setlength{\parindent}{0pt}
\setlength{\parskip}{0.35em plus 0.10em minus 0.08em}
\setlength{\emergencystretch}{3em}
\setlist{itemsep=0.25em, topsep=0.40em, parsep=0pt}
\allowdisplaybreaks[2]
% === 错题本排版优化结束 ===
\begin{document}
\section*{华东理工大学2026年数学分析真题}
\section*{54．华东理工大学2026年数学分析试卷}

\subsection{题目1}
1．写出基本数列（Cauchy 数列）的定义。

\subsection{题目2}
2．设 \(x_{1}>0\) ，对 \(n \geq 1\) ，有 \(x_{n+1}=\arctan x_{n}\) ，证明：\(\left\{x_{n}\right\}\) 是基本数列．

二．设函数 \(f\) 在闭区间 \([a, b]\) 上有二阶导数，且 \(f^{\prime}(a)=f^{\prime}(b)=0\) ，证明：存在 \(\xi \in(a, b)\) ，使得

\[
\left|f^{\prime \prime}(\xi)\right| \geq 4 \frac{|f(b)-f(a)|}{(b-a)^{2}} .
\]

三．在 \(\mathbb{R}^{2}\) 上的函数

\[
f(x, y)= \begin{cases}\frac{x^{2}-y^{2}}{x^{2}+y^{2}}, & x^{2}+y^{2} \neq 0 \\ 0, & x^{2}+y^{2}=0\end{cases}
\]

讨论 \(f\) 在 \((0,0)\) 点的连续性．

\subsection{题目3}
1．\(f\) 在 \([0,1]\) 上可积，说明： \(\int_{0}^{1} f(x) \mathrm{d} x=\lim _{n \rightarrow \infty} \sum_{k=1}^{n} \frac{1}{n} f\left(\frac{k-1}{n}\right)\) ．

\subsection{题目4}
2．对 \(n \geq 1\) ，定义 \(A_{n}=\frac{\sum_{k=1}^{n} k \sin \left(\frac{k-1}{n} \pi\right)}{n^{2}}\) ，求 \(\lim _{n \rightarrow \infty} A_{n}\) ．

\begin{itemize}
  \item[五．] 记 \(I_{p}=\int_{2}^{+\infty} \frac{(\ln x)^{p} \sin x}{x} \mathrm{~d} x, p \in \mathbb{R}\) ．

  \item[1．] 当 \(p=0\) 时，求 \(I_{p}\) 的敛散性。

  \item[2．] 当 \(p=1\) 时，求 \(I_{p}\) 的敛散性。
\end{itemize}

六．计算曲线积分 \(I=\oint_{\Gamma}|\sqrt{3} y-x| \mathrm{d} x-5 z \mathrm{~d} z\) ，其中\\
\(\Gamma:\left\{\begin{array}{l}x^{2}+y^{2}+z^{2}=8 \\ x^{2}+y^{2}=2 z\end{array}\right.\)\\
从 \(z\) 轴正向往坐标原点看去取逆时针方向。

七．已知 \(F\) 有二阶连续偏导数，\(z=z(x, y)\) 是由 \(F\left(z+\frac{1}{x}, z-\frac{1}{y}\right)=0\) 确定的隐函数．且 \(z\) 有连续二阶偏导数。证明：

\begin{itemize}

  \item[1．] \(x^{2} \frac{\partial z}{\partial x}-y^{2} \frac{\partial z}{\partial y}=1\) ．

  \item[2．] \(x^{3} \frac{\partial^{2} z}{\partial x^{2}}+x y(x-y) \frac{\partial^{2} z}{\partial x \partial y}-y^{3} \frac{\partial^{2} z}{\partial y^{2}}+2=0\) ．
\end{itemize}

八．设 \(\sum\) 是单位球面 \(\left\{(x, y, z) \in \mathbb{R}^{3}: x^{2}+y^{2}+z^{2}=1\right\}\) ，记 \(I_{p}=\iint_{\sum}(x+y+z)^{p} \mathrm{~d} S, p\) 是正整数，求\(I_{1}, I_{2}\) 。

\subsection{题目9}
1．给出函数列一致收敛的定义。

\subsection{题目10}
2．对 \(n \geq 1\) ，设 \(S_{n}(x)=\frac{x}{1+n^{2} x^{2}}\) ，证明：函数列 \(\left\{S_{n}(x)\right\}\) 在 \((0,+\infty)\) 上一致收敛。

\subsection{题目11}
3．当 \(n \geq 1\) 时，\(T_{n}(x)=\frac{\sqrt{x}}{1+n^{2} x^{2}}\) ，问：函数列 \(\left\{T_{n}(x)\right\}\) 在 \((0,+\infty)\) 上是否一致收敛？请证明你的结论．
\end{document}
